% Part:sets-functions-relations
% Chapter: sets
% Section: reduction

\documentclass[../../../include/open-logic-section]{subfiles}

\begin{document}

\olfileid{sfr}{siz}{red}

\olsection{ન્યૂનીકરણ}

\begin{editorial}
  આ વિભાગ \olref[nen]{sec}નાં પરિણામોને અનુરૂપ રીતે
  ન્યૂનીકરણ દ્વારા અગણનીયતા સાબિત કરે છે. \olref[nen-alt]{sec}નાં
  પરિણામોને અનુરૂપ થોડું વધુ સંક્ષિપ્ત વૈકલ્પિક સ્વરૂપ
  \olref[red-alt]{sec}માં આપેલું છે.
\end{editorial}

વિકર્ણીકરણની દલીલ વડે આપણે $\Pow{\PosInt}$ એ !!{nonenumerable}
છે તે બતાવ્યું. $0$ અને $1$ના બધા અનંત અનુક્રમોનો ગણ
$\Bin^\omega$ !!{nonenumerable} છે તેની સાબિતી આપણી પાસે
પહેલેથી હતી. $\Pow{\PosInt}$ !!{nonenumerable} છે તે સાબિત કરવાની
આ બીજી રીત છે: \emph{જો $\Pow{\PosInt}$ એ !!{enumerable} હોય, તો
$\Bin^\omega$ પણ !!{enumerable} છે} તે બતાવો. $\Bin^\omega$
!!{enumerable} નથી તે આપણે જાણીએ છીએ, તેથી $\Pow{\PosInt}$ પણ
હોઈ શકે નહીં. તેને એક સમસ્યાનું બીજી સમસ્યામાં \emph{ન્યૂનીકરણ}
કહે છે---અહીં આપણે $\Bin^\omega$ને પરિગણિત કરવાની સમસ્યાનું
$\Pow{\PosInt}$ને પરિગણિત કરવાની સમસ્યામાં ન્યૂનીકરણ કરીએ છીએ.
બીજી સમસ્યાનો ઉકેલ---$\Pow{\PosInt}$ની પરિગણના---પહેલી સમસ્યાનો
ઉકેલ---$\Bin^\omega$ની પરિગણના---આપશે.

ગણ~$B$ને પરિગણિત કરવાની સમસ્યાનું ગણ~$A$ને પરિગણિત કરવાની
સમસ્યામાં ન્યૂનીકરણ કેવી રીતે કરવું? $A$ની પરિગણનાને $B$ની
પરિગણનામાં ફેરવવાની રીત આપીએ. એ કરવાની સૌથી સરળ રીત
!!a{surjective} વિધેય $f\colon A \to B$ વ્યાખ્યાયિત કરવાની છે.
જો $x_1$, $x_2$, \dots{} એ $A$ની પરિગણના કરે, તો
$f(x_1)$, $f(x_2)$, \dots{} એ $B$ની પરિગણના કરશે.
આપણા કિસ્સામાં આપણે વ્યાપ્ત વિધેય
$f\colon \Pow{\PosInt} \to \Bin^\omega$ શોધીએ છીએ.

\begin{prob}
બતાવો કે જો કોઈ !!{injective} વિધેય $g\colon B \to A$ હોય અને
$B$ એ !!{nonenumerable} હોય, તો $A$ પણ અગણનીય છે.
$A$ની પરિગણનામાંથી $B$ની પરિગણના બનાવવા $g$નો ઉપયોગ કેવી રીતે
કરી શકાય તે બતાવીને આ કરો.
\end{prob}

\begin{proof}[ન્યૂનીકરણ દ્વારા {\olref[nen]{thm:nonenum-pownat}}ની સાબિતી]
ધારો કે $\Pow{\PosInt}$ !!{enumerable} હોત અને તેથી તેની કોઈ
પરિગણના $Z_{1}$, $Z_{2}$, $Z_{3}$, \dots હોત.

વિધેય $f \colon \Pow{\PosInt} \to \Bin^\omega$ને આ રીતે
વ્યાખ્યાયિત કરો: $f(Z)$ એ એવો અનુક્રમ~$s$ છે કે
$s(n) = 1$ ત્યારે અને માત્ર ત્યારે જ જ્યારે $n \in Z$, અને
અન્યથા $s(n) = 0$. આ સ્પષ્ટપણે વિધેય વ્યાખ્યાયિત કરે છે,
કારણ કે જ્યારે પણ $Z \subseteq \PosInt$ હોય ત્યારે કોઈ પણ
$n \in \PosInt$ એ $Z$નો !!a{element} હોય અથવા ન હોય.
દાખલા તરીકે, ધન યુગ્મ સંખ્યાઓનો ગણ $2\PosInt = \{2,
4, 6, \dots\}$ એ અનુક્રમ $010101\dots$ પર મોકલાય છે,
ખાલી ગણ $0000\dots$ પર અને $\PosInt$ પોતે $1111\dots$ પર મોકલાય છે.
\sourcecorrection{OLSIZ-004}{મૂળની \texttt{reduction.tex}, પંક્તિઓ
51--58માં \(f(Z)\)ની કિંમત માટે અનિર્ધારિત સૂચકવાળું \(s_k\) લખાયું
છે, જ્યારે પછીની વ્યાપ્તતાની દલીલ \(s\) વાપરે છે. અહીં આખી
વ્યાખ્યામાં \(s\) સુસંગત રીતે વાપર્યું છે.}

તે !!{surjective} પણ છે: $0$ અને $1$નો દરેક અનુક્રમ ધન
પૂર્ણાંકોના કોઈ ગણને અનુરૂપ છે, એટલે કે અનુક્રમમાં જે સ્થાનો પર
$1$ હોય તેને અનુરૂપ પૂર્ણાંકો તેના સભ્યો છે. વધુ ચોક્કસ રીતે,
ધારો કે $s \in \Bin^\omega$. આ રીતે $Z \subseteq \PosInt$
વ્યાખ્યાયિત કરો:
\[
Z = \Setabs{n \in \PosInt}{s(n) = 1}
\]
ત્યારે $f$ની વ્યાખ્યા તપાસતાં $f(Z) = s$ મળે છે.

હવે આ યાદીનો વિચાર કરો:
\[
f(Z_1), f(Z_2), f(Z_3), \dots
\]
$f$ એ !!{surjective} હોવાથી $\Bin^\omega$નો દરેક સભ્ય કોઈ
દલીલ માટેની $f$ની કિંમત તરીકે આવવો જોઈએ અને તેથી યાદીમાં આવવો જ
જોઈએ. આ યાદી આથી આખા~$\Bin^\omega$ની પરિગણના કરે છે.

એટલે જો $\Pow{\PosInt}$ એ !!{enumerable} હોત, તો $\Bin^\omega$
!!{enumerable} હોત. પરંતુ $\Bin^\omega$ એ !!{nonenumerable} છે
(\olref[nen]{thm:nonenum-bin-omega}). તેથી $\Pow{\PosInt}$ એ
!!{nonenumerable} છે.
\end{proof}

\begin{explain}
ન્યૂનીકરણ કઈ દિશામાં જાય છે તેમાં ગૂંચવાવું સહેલું છે.
દાખલા તરીકે, !!a{surjective} વિધેય $g \colon \Bin^\omega \to B$
હોવાથી $B$ એ !!{nonenumerable} છે એવું \emph{સ્થાપિત થતું નથી}.
($g(s) = s(1)$ વડે વ્યાખ્યાયિત વિધેય
$g \colon \Bin^\omega \to \Bin$નો વિચાર કરો: તે $0$ અને $1$ના
અનુક્રમને તેના પહેલા !!{element} પર મોકલે છે. તે !!{surjective} છે,
કારણ કે કેટલાક અનુક્રમો $0$થી અને કેટલાક $1$થી શરૂ થાય છે.
પરંતુ $\Bin$ સાન્ત છે.) વળી, વિધેય~$f$ !!{surjective} હોવું જ
જોઈએ; અન્યથા દલીલ ચાલતી નથી: $f(x_1)$, $f(x_2)$, \dots{}માં
$B$ના બધા !!{element}s આવે તેની ખાતરી રહેતી નથી. દાખલા તરીકે,
\[
h(n) = \underbrace{000\dots0}_{\text{$n$ $0$'s}}111\dots
\]
વિધેય $h\colon \PosInt \to \Bin^\omega$ વ્યાખ્યાયિત કરે છે,
પરંતુ $\PosInt$ એ !!{enumerable} છે.
\sourcecorrection{OLSIZ-005}{મૂળની \texttt{reduction.tex}, પંક્તિઓ
85--101માં \(h(n)\) માટે માત્ર \(n\) શૂન્યોની સાન્ત પ્રતીકશ્રેણી
આપેલી છે, જે \(\Bin^\omega\)ની કિંમત નથી. અહીં તે પછી અનંત ઘણા
એકડા ઉમેરીને અનંત દ્વિઅંકી અનુક્રમ બનાવ્યો છે; દલીલને માત્ર
આવા કોઈ અવ્યાપ્ત વિધેયની જરૂર છે.}
\end{explain}

\begin{prob}
ન્યૂનીકરણની દલીલ વડે બતાવો કે ધન પૂર્ણાંકોની જોડોના બધા
\emph{ગણોનો} ગણ !!{nonenumerable} છે.
\end{prob}

\begin{prob}\label{sfr:siz:red:prob:nat-nat}
  ન્યૂનીકરણની દલીલ વડે બતાવો કે બધા વિધેયો
  $f\colon \Nat \to \Nat$નો ગણ~$X$ !!{nonenumerable} છે.
  (સંકેત: $X$થી $\Bin^\omega$માં વ્યાપ્ત વિધેય આપો.)
\end{prob}

\begin{prob}
ન્યૂનીકરણની દલીલ વડે બતાવો કે પ્રાકૃતિક સંખ્યાઓના અનંત
અનુક્રમોનો ગણ~$\Nat^\omega$ !!{nonenumerable} છે.
\end{prob}

\begin{prob}
$P$ એ ધન પૂર્ણાંકોના ગણથી ગણ~$\{0\}$માંના વિધેયોનો ગણ અને
$Q$ એ ધન પૂર્ણાંકોના ગણથી ગણ~$\{0\}$માંના \emph{આંશિક}
વિધેયોનો ગણ હોય. બતાવો કે $P$ એ !!{enumerable} છે અને $Q$ નથી.
(સંકેત: $\Bin^\omega$ને પરિગણિત કરવાની સમસ્યાનું $Q$ને
પરિગણિત કરવાની સમસ્યામાં ન્યૂનીકરણ કરો.)
\end{prob}

\begin{prob}
$S$ એ ધન પૂર્ણાંકોના ગણથી ગણ \{0,1\}માંના બધા
!!{surjective} વિધેયોનો ગણ હોય; એટલે કે $S$માં બધા
!!{surjective}~$f\colon \PosInt \to \Bin$ છે. બતાવો કે $S$ એ
!!{nonenumerable} છે.
\end{prob}

\begin{prob}
બતાવો કે બધી વાસ્તવિક સંખ્યાઓનો ગણ~$\Real$ !!{nonenumerable} છે.
\end{prob}

\end{document}
